\chapter{Предиктивное кодирование, сопоставление представлений и адаптивное вычисление}
\label{ch:related-work}

Эта глава задаёт литературный контекст для трёх исследовательских вопросов
диссертации. Сначала рассматривается обратное распространение ошибки как
базовая вычислительная линия обучения глубоких сетей. Затем обсуждается переход
от нейробиологической постановки предиктивного кодирования (PC) к современным
вариантам обучения через вывод и к вопросу о связи их обновлений с обратным
распространением ошибки (BP). Отдельно рассматриваются методы сопоставления
внутренних представлений и градиентов. Наконец, литература по адаптивным
вычислениям, раннему выходу и селективному предсказанию используется для
позиционирования QWake как задачи безопасного решения о продолжении или
прекращении вычисления.

Важно различать несколько типов утверждений, которые в литературе иногда
оказываются рядом, но не являются взаимозаменяемыми: сходство итогового
качества, близость внутренних представлений, близость направлений градиентов,
точное совпадение обновлений при специальных условиях и вычислительная
эффективность. Настоящая диссертация сохраняет эти уровни раздельными и
использует каждый источник только в пределах того вопроса, который он
непосредственно рассматривает.

\section{Обратное распространение ошибки как базовая вычислительная линия}
\label{sec:related-bp}

Классическая работа Rumelhart, Hinton и Williams формализовала практическое
обучение многослойных сетей посредством обратного распространения ошибки и
связала изменение скрытых представлений с минимизацией ошибки на выходе
\cite{rumelhart1986backprop}. В современной терминологии BP можно рассматривать
как эффективную реализацию цепного правила для вычисления производных по
параметрам композиции функций. Для данной диссертации BP важен не как
нейробиологическая модель, а как хорошо определённый эталонный алгоритмический
режим, относительно которого можно строить парные сравнения.

Такое позиционирование существенно для интерпретации результатов. Совпадение
или близость метрики качества PC и BP не означает совпадения внутренних
траекторий обучения. Аналогично, близость градиентов одного слоя не переносится
автоматически на остальные слои, а приближение к BP при специальных условиях
не определяет стоимость достижения этих условий. Поэтому в экспериментальной
части BP используется одновременно как поведенческая базовая линия, как
источник эталонных градиентов и как точка отсчёта для сопоставленного
профилирования.

\section{Иерархическое предиктивное кодирование}
\label{sec:related-pc}

Историческая линия предиктивного кодирования восходит к моделям, в которых
соседние уровни иерархии обмениваются предсказаниями и ошибками предсказания.
Rao и Ballard связали такую архитектуру с обработкой зрительной информации и
иерархическим объяснением входа \cite{rao1999predictive}. В этой постановке
единицы ошибки предсказания имеют самостоятельный вычислительный смысл, а
состояние сети определяется не только параметрами, но и внутренними
переменными, изменяющимися в ходе итеративного вывода.

Для машинного обучения особенно важна работа Whittington и Bogacz
\cite{whittington2017approximation}. В ней обучение с учителем в PC-сети с
локальными хеббовскими обновлениями сопоставляется с BP; при определённых
параметрах обновления PC приближаются к обновлениям обратного распространения
ошибки. Эта линия переводит вопрос из общего сравнения нейробиологических
принципов в более узкую алгоритмическую задачу: при каких динамиках, начальных
условиях и правилах обновления локальная процедура вывода воспроизводит или
приближает эталонный градиент.

Современные работы также рассматривают способы ускорения и изменения внутренней
динамики вывода. Например, Hybrid Predictive Coding сочетает несколько
временных режимов вывода, мотивируя разделение быстрых и более медленных
вычислительных процессов \cite{tschantz2023hybrid}. Для настоящей диссертации
эта линия важна как контекст: итеративность PC делает число внутренних
вычислительных шагов самостоятельной величиной, а следовательно позволяет
отдельно исследовать достаточность промежуточного состояния и цену его
дальнейшего уточнения.

\section{Связь предиктивного кодирования и обратного распространения ошибки}
\label{sec:related-pc-bp}

Связь PC и BP не сводится к одному общему утверждению об эквивалентности.
Разные работы используют разные варианты динамики, инициализации и порядка
обновлений, поэтому область применимости результата должна сохраняться вместе
с самим результатом.

Whittington и Bogacz рассматривают приближение BP локальными правилами
\cite{whittington2017approximation}. Song и соавторы вводят условия, при которых
вариант сети предиктивного кодирования воспроизводит обновления BP точно
\cite{song2020exact}. Millidge, Tschantz и Buckley расширяют анализ на
произвольные дифференцируемые вычислительные графы и описывают асимптотическое
схождение градиентов PC к градиентам обратного распространения ошибки в
рассматриваемой конструкции \cite{millidge2022arbitrary}.

Эти результаты не устраняют вопрос о том, какие именно ограничения необходимы
для сопоставления. Rosenbaum анализирует отношения между PC и BP и разделяет
несколько режимов, в которых выводы об эквивалентности имеют различную силу
\cite{rosenbaum2022relationship}. Для этой работы учитывается также официальная
поправка 2025 года \cite{rosenbaum2025correction}; поэтому литературная
интерпретация исходной статьи не используется без учёта исправленной версии.

Отдельную линию составляет вычислительная критика точных или близких к BP
вариантов PC. Zahid, Guo и Fountas рассматривают современные варианты PC как
возможную нейроморфную альтернативу и выводят ограничения по временной
сложности для исследуемых конструкций \cite{zahid2023critical}. Эта работа
особенно релевантна RQ2: локальность обновлений или теоретическая связь с BP
сама по себе не устанавливает вычислительное преимущество. Поэтому в
диссертации отнесение наблюдаемого эффекта к механизму и сопоставленное
профилирование стоимости рассматриваются как отдельная эмпирическая задача
после проверки поведенческой и градиентной сопоставимости.

Таким образом, литературный контекст поддерживает последовательную структуру
настоящей работы: сначала требуется установить, \emph{что именно} совпадает или
различается между режимами, затем локализовать механизм расхождения и только
после этого обсуждать вычислительную цену конкретной реализации.

\section{Сходство внутренних представлений и градиентов}
\label{sec:related-representations}

Сравнение нейронных сетей только по точности или функции потерь оставляет
открытым вопрос, формируются ли внутри модели сходные представления. SVCCA
объединяет сингулярное разложение и канонический корреляционный анализ для
сопоставления активаций разных сетей или разных моментов обучения
\cite{raghu2017svcca}. Morcos, Raghu и Bengio развивают анализ представлений на
основе CCA и обсуждают устойчивость выводов о сходстве
\cite{morcos2018insights}.

Kornblith и соавторы систематически сравнивают семейство мер сходства
представлений и вводят центрированное выравнивание ядер (CKA) как удобную
метрику сопоставления представлений \cite{kornblith2019similarity}. В
диссертации CKA используется как основная диагностика на уровне
представлений, тогда как методы на основе CCA задают дополнительный
литературный контекст.

Сходство представлений и сходство градиентов отвечают на разные вопросы. Первое
описывает геометрию активаций на выбранной выборке, второе --- локальное
направление обновления параметров. Поэтому этап~3A не агрегирует эти измерения
в единую ``оценку сходства''. Косинусное сходство градиентов, отношение их норм
и CKA рассматриваются как взаимодополняющие координаты внутреннего состояния
обучения. Такое разделение особенно важно для режимов, которые могут сохранять
сходную геометрию представлений при заметно иной величине градиентного сигнала
в ранних слоях.

\section{Адаптивные вычисления и раннее прекращение вычисления}
\label{sec:related-adaptive-computation}

Идея делать объём вычисления зависимым от конкретного входа существует
независимо от предиктивного кодирования. Adaptive Computation Time (ACT)
позволяет рекуррентной сети адаптировать число внутренних шагов перед выдачей
результата \cite{graves2016adaptive}. Это формулирует вычислительный бюджет как
часть решения модели, а не как фиксированную константу архитектуры.

Для сетей прямого распространения раннее завершение часто реализуется с помощью
промежуточных выходов. Bolukbasi и соавторы формулируют адаптивную оценку сети
как выбор того, какие компоненты сети необходимо вычислять для конкретного
примера, и обучают правило раннего выхода с явной целью сокращения
вычислительного времени \cite{bolukbasi2017adaptive}. В более поздней литературе
архитектуры с несколькими выходами остаются активной областью исследования:
например, Kubaty и соавторы анализируют, как способ совместного или раздельного
обучения основной сети и выходных голов влияет на свойства модели с несколькими
выходами \cite{kubaty2025multiexit}.

Эта литература близка QWake по общей экономической постановке: дополнительный
механизм решения должен окупаться за счёт предотвращённого вычисления. Однако
объект решения различается. Типичный классификатор с ранним выходом добавляет
промежуточную предсказательную головку. QWake в настоящей работе не вводит новую
классификационную головку для каждого вычислительного шага; вместо этого
анализируется, можно ли по уже доступному состоянию до действия безопасно
распознать состояние, после которого продолжение зарегистрированного вычисления
PC не требуется в пределах зафиксированного критерия.

\section{Селективное предсказание, отказ от ответа и безопасный выход}
\label{sec:related-selective}

Ранний выход создаёт второй вопрос помимо экономии: когда сокращённое вычисление
допустимо с точки зрения качества результата. В селективном предсказании модель
может отказаться от ответа для части входов и тем самым формирует компромисс
между риском и покрытием. SelectiveNet объединяет функцию выбора и
предсказательную модель в одну обучаемую архитектуру
\cite{geifman2019selectivenet}. Для диссертации существенен сам принцип
разделения покрытия и риска: высокая доля принимаемых состояний не имеет смысла,
если среди них присутствуют состояния, которые критерий считает
недостаточными.

Современные работы по раннему выходу делают аспект безопасности явным. Jazbec и
соавторы применяют схему управления риском к сетям с ранним выходом, связывая
решение о выходе с заранее заданной границей качества
\cite{jazbec2024riskcontrol}. Другая работа той же линии рассматривает
последовательную согласованность множеств неопределённости между выходами сети с
несколькими выходами \cite{jazbec2024nested}. В 2025 году EERO объединяет ранний
выход и отказ от ответа с явным ограничением вычислительного бюджета
\cite{valade2025eero}.

QWake находится рядом с этой литературой, но использует более жёсткий
ограниченный критерий безопасности. На калибровочных доказательных материалах
правило считается допустимым только при отсутствии опасных принятых решений;
ненулевого покрытия самого по себе недостаточно. После проверки безопасности
применяется экономический критерий по совокупной чистой экономии. Поэтому QWake
можно рассматривать как специальный случай селективного вычисления: принятие
означает признание текущего состояния достаточным и прекращение оставшегося
вычисления, а отказ --- продолжение канонического вычисления по принципу
безопасного отказа. Такое сопоставление является концептуальным; QWake не
заимствует статистические гарантии методов управления риском или конформных
методов, если они не реализованы и не проверены непосредственно в
экспериментальном протоколе.

\section{Позиционирование настоящей работы}
\label{sec:related-positioning}

Литература задаёт несколько отдельных осей, которые настоящая диссертация
соединяет в одной причинной экспериментальной линии:

\begin{enumerate}
  \item PC может приближать или при специальных условиях воспроизводить
        определённые обновления BP, но область такого сопоставления зависит от
        конкретного алгоритмического режима
        \cite{whittington2017approximation,song2020exact,millidge2022arbitrary,
        rosenbaum2022relationship,rosenbaum2025correction};

  \item сходство на уровне представлений требует отдельной диагностики и не
        следует автоматически из близости итоговой ошибки
        \cite{raghu2017svcca,morcos2018insights,kornblith2019similarity};

  \item вычислительная цена механизма должна измеряться отдельно от его
        функциональной связи с BP
        \cite{zahid2023critical};

  \item адаптивные вычисления создают компромисс между полнотой вычисления и
        стоимостью, а литература по селективному предсказанию и раннему выходу
        добавляет явный аспект соотношения безопасности и покрытия
        \cite{graves2016adaptive,bolukbasi2017adaptive,
        geifman2019selectivenet,jazbec2024riskcontrol,valade2025eero}.
\end{enumerate}

Отсюда следует позиция диссертации относительно предыдущих работ. Её вклад не
формулируется как ещё одно общее утверждение об эквивалентности PC и BP и не
сводится к построению новой архитектуры с несколькими выходами. Исследование
последовательно проверяет поведение, внутренние представления и градиенты,
механизм и измеренную стоимость в одной версионируемой экспериментальной линии,
а затем задаёт отдельный ограниченный вопрос о возможности безопасно прекратить
уже существующее итеративное вычисление. Именно это разделение мотивирует
RQ1--RQ3 и структуру последующих глав.
